\hypertarget{python-3.13-copy-and-patch-jit-compiler-architecture}{%
\chapter{Python 3.13: Copy-and-Patch JIT Compiler
Architecture}\label{python-3.13-copy-and-patch-jit-compiler-architecture}}

\hypertarget{copy-and-patch-jit-compiler-pipeline}{%
\subsection{21.1 Copy-and-Patch JIT Compiler
Pipeline}\label{copy-and-patch-jit-compiler-pipeline}}

Python 3.13 introduces an experimental \textbf{Copy-and-Patch
Just-In-Time (JIT) compiler} (PEP 744). Unlike classic template JITs or
heavy tracing compilers (like PyPy or V8), a copy-and-patch JIT achieves
fast compilation times with very low runtime memory overhead.

\hypertarget{compile-time-vs.-runtime-pipeline}{%
\subsubsection{1. Compile-Time vs.~Runtime
Pipeline}\label{compile-time-vs.-runtime-pipeline}}

\begin{lstlisting}
BUILD TIME (Clang/LLVM):
[CPython C Code] -> [LLVM Clang] -> [Machine Code Stencils (ELF)] -> [Stencil Assembly Extractor] -> [Stitcher Headers]

RUNTIME EXECUTION:
[Bytecode Instructions] -> [JIT Engine] -> [Copy Stencils to Memory] -> [Patch Relocations/Variables] -> [Execute Native Code]
\end{lstlisting}

\begin{enumerate}
\def\labelenumi{\arabic{enumi}.}
\tightlist
\item
  \textbf{Build Time}: The CPython build system uses LLVM/Clang to
  compile individual bytecode execution paths (helpers in
  \passthrough{\lstinline!ceval.c!}) into object file templates
  (stencils). A python script parses these ELF files and extracts the
  raw machine instructions, noting the offsets of ``holes''
  (relocations) representing variables, stack offsets, or jump
  addresses.
\item
  \textbf{Runtime}: When a bytecode sequence becomes ``hot,'' the JIT
  engine copies the pre-compiled stencils into an executable memory page
  and patches the holes with runtime addresses (e.g., target object
  references or jump offsets).
\end{enumerate}

\begin{center}\rule{0.5\linewidth}{0.5pt}\end{center}

\hypertarget{stencil-relocation-and-hole-patching}{%
\subsection{21.2 Stencil Relocation and Hole
Patching}\label{stencil-relocation-and-hole-patching}}

Consider a bytecode operation like \passthrough{\lstinline!LOAD\_FAST!}
inside a stencil template:

\begin{lstlisting}
mov rax, [r14 + HOLE_OFFSET_LOCAL_VAR]  ; Load local variable offset
\end{lstlisting}

During the patching phase, the JIT engine replaces the placeholder
\passthrough{\lstinline!HOLE\_OFFSET\_LOCAL\_VAR!} with the actual frame
offset:
\[\text{Relocation Address} = \text{Frame Base} + (\text{Var Index} \times 8)\]
The engine then calls \passthrough{\lstinline!mprotect()!} to set the
memory page's permissions to Read/Execute (RX), enabling direct
hardware-native execution of the compiled sequence.

\begin{center}\rule{0.5\linewidth}{0.5pt}\end{center}

\hypertarget{jit-performance-memory-footprint-bounds}{%
\subsection{21.3 JIT Performance \& Memory Footprint
Bounds}\label{jit-performance-memory-footprint-bounds}}

Because the JIT engine only performs simple memory copying
(\passthrough{\lstinline!memcpy!}) and value patching, it avoids the
heavy parsing, graph building, register allocation, and optimization
passes associated with traditional JIT compilers. * \textbf{Compilation
Cost}: Compiling code is extremely fast, requiring only microseconds. *
\textbf{Memory Footprint}: Minimal overhead, as the template stencils
are compiled during CPython's build time and stored in the static binary
data section. * \textbf{Instruction Cache Friendly}: Generated machine
code templates match the specialized adaptive bytecodes from Chapter 17,
ensuring high hardware pipeline utilization.

\begin{center}\rule{0.5\linewidth}{0.5pt}\end{center}
